% GENERATED FILE. Edit canonical lessons, then run npm run generate:books.
% canonical-source-hash: fnv1a64:c35b3dcceeea407e
% canonical-lessons: ZH-C78-duoda, ZH-C78-zhangda, ZH-C78-qian, ZH-W78-qian-before, ZH-C78-nianji, ZH-C78-xiaoshihou

\chapter{Age and Growing Up}
\label{ch:zh-age-and-growing-up}
\hlchaptermodality{78}

\begin{chapteropening}
\textbf{By the end of this chapter:} \emph{I can ask how old somebody is and talk about growing up and years ago.}

Complete the last lesson of chapter 78: I can ask how old somebody is and talk about growing up and years ago.
\end{chapteropening}

\section[\texorpdfstring{\zh{多大} (duō dà)}{多大 (duō dà)}]{\zh{多大} (duō dà) — how old (asking a person's age)}

\label{lesson:ZH-C78-duoda}



\begin{quote}
Before the new one: say the Mandarin for extremely, then the Mandarin for really.
\end{quote}

\subsection*{You'll want to know: \zh{多大}}
\textbf{\zh{多大}} — \emph{duō dà} — \textquotedblleft{}how old (asking a person's age)\textquotedblright{}.

\begin{quote}
\textbf{\zh{生日快乐你今年多大}} \emph{shēngrì kuàilè nǐ jīnnián duō dà} — Happy birthday! How old are you now?
\end{quote}

\begin{cousinweb}
\textbf{\zh{多大}} (\emph{duō dà}) is built from \zh{多} 'much' and \zh{大} 'big': how big, said of a person's years.
\end{cousinweb}

\begin{grammarlens}[title={What you've built}]
One more word about people.
\end{grammarlens}

\subsection*{Your turn}
\begin{itemize}
\raggedright
  \item \emph{Say it:} \emph{duō dà}
  \item \emph{Say it:} \emph{duō dà}, once more
  \item \emph{Say it:} \emph{shēngrì kuàilè nǐ jīnnián duō dà} — Happy birthday! How old are you now?
  \item \emph{Recall:} say the Mandarin for extremely, then the Mandarin for really, then say \emph{duō dà} again
\end{itemize}

\begin{tcolorbox}[breakable,colback=teal!4,colframe=teal!35!black,title={Before you move on}]
What does \zh{多大} mean? (\textquotedblleft{}How old (asking a person's age)\textquotedblright{}.) Say it once more.
\end{tcolorbox}
\section[\texorpdfstring{\zh{长大} (zhǎngdà)}{长大 (zhǎngdà)}]{\zh{长大} (zhǎngdà) — to grow up}

\label{lesson:ZH-C78-zhangda}



\begin{quote}
Before the new one: say the Mandarin for really, then the Mandarin for how old.
\end{quote}

\subsection*{You'll want to know: \zh{长大}}
\textbf{\zh{长大}} — \emph{zhǎngdà} — \textquotedblleft{}to grow up\textquotedblright{}.

\begin{quote}
\textbf{\zh{他们长大了}} \emph{tāmen zhǎngdà le} — They've grown up.
\end{quote}

The character you wrote for cháng, long, is read zhǎng here, to grow.

\begin{cousinweb}
\textbf{\zh{长大}} (\emph{zhǎngdà}) is built from zhǎng 'to grow' and \zh{大} 'big', both inherited from Middle Chinese, \emph{trjangX} and \emph{dajH} (Guangyun, in Baxter's transcription). The dictionary glosses the first as 'big' and 'elder'; growing is not given there, so how it got that sense is not settled.
\end{cousinweb}

\begin{grammarlens}[title={What you've built}]
One more word about people.
\end{grammarlens}

\subsection*{Your turn}
\begin{itemize}
\raggedright
  \item \emph{Say it:} \emph{zhǎngdà}
  \item \emph{Say it:} \emph{zhǎngdà}, once more
  \item \emph{Say it:} \emph{tāmen zhǎngdà le} — They've grown up.
  \item \emph{Recall:} say the Mandarin for really, then the Mandarin for how old, then say \emph{zhǎngdà} again
\end{itemize}

\begin{tcolorbox}[breakable,colback=teal!4,colframe=teal!35!black,title={Before you move on}]
What does \zh{长大} mean? (\textquotedblleft{}To grow up\textquotedblright{}.) Say it once more.
\end{tcolorbox}
\section[\texorpdfstring{\zh{前} (qián)}{前 (qián)}]{\zh{前} (qián) — before; ago; front}

\label{lesson:ZH-C78-qian}



\begin{quote}
Before the new one: say the Mandarin for how old, then the Mandarin for to grow up.
\end{quote}

\subsection*{You'll want to know: \zh{前}}
\textbf{\zh{前}} — \emph{qián} — \textquotedblleft{}before; ago; front\textquotedblright{}.

\begin{quote}
\textbf{\zh{三年前我八岁}} \emph{sān nián qián wǒ bā suì} — Three years ago I was eight.
\end{quote}

The new character, in pieces: \zh{前} — two horns over a bar, the moon piece read yuè on the left, and a blade of two verticals on the right. The next lesson puts it in your hand.

\begin{cousinweb}
\textbf{\zh{前}} (\emph{qián}) is inherited from Middle Chinese \emph{dzen} (Guangyun, in Baxter's transcription), glossed 'before, first'.
\end{cousinweb}

\begin{grammarlens}[title={What you've built}]
One more word about people.
\end{grammarlens}

\subsection*{Your turn}
\begin{itemize}
\raggedright
  \item \emph{Say it:} \emph{qián}
  \item \emph{Say it:} \emph{qián}, once more
  \item \emph{Say it:} \emph{sān nián qián wǒ bā suì} — Three years ago I was eight.
  \item \emph{Recall:} say the Mandarin for how old, then the Mandarin for to grow up, then say \emph{qián} again
\end{itemize}

\begin{tcolorbox}[breakable,colback=teal!4,colframe=teal!35!black,title={Before you move on}]
What does \zh{前} mean? (\textquotedblleft{}Before; ago; front\textquotedblright{}.) Say it once more.
\end{tcolorbox}
\section[\texorpdfstring{\zh{前} (qián)}{前 (qián)}]{\zh{前} — horns, a bar, the moon and a blade}

\label{lesson:ZH-W78-qian-before}



\begin{quote}
The word first, then the character inside it.

\begin{itemize}
\raggedright
  \item \emph{Recall:} say \emph{qián} — \textbf{R1}, one lesson back, before it decays
\end{itemize}

\begin{quote}
\zh{三} \zh{年}
\end{quote}

\emph{Read it:} each one aloud, with nothing said first — \emph{sān}, \emph{nián}
\end{quote}

\subsection*{Script — horns, a bar, the moon and a blade}
\hlblockfigure{\detokenize{figures/ZH-W78-qian-before-filmstrip.pdf}}{How \zh{前} is written, stroke by stroke}

\begin{quote}
\zh{前}
\end{quote}

\textbf{Two horns, two strokes}: a dot on the left, a short left-falling stroke on the right.

\textbf{A long horizontal} under them.

\textbf{The moon piece, four strokes, read yuè} as in the word for month: a left side that falls down, a top that turns down the right side and hooks, and two short bars inside.

\textbf{The blade, two strokes}: a short vertical, then a long vertical that hooks to the left.

2 + 1 + 4 + 2 = nine.

\subsection*{Your turn}
\begin{itemize}
\raggedright
  \item \emph{Write it:} \zh{前}, all nine, with the model in front of you
  \item \emph{Cover:} write it again from memory, and count the pieces, not the strokes
\end{itemize}

\begin{tcolorbox}[breakable,colback=teal!4,colframe=teal!35!black,title={Before you move on}]
How many strokes in \zh{前}? (\textbf{Nine}.) What comes after the horns? (\textbf{The long bar}.) Which stroke hooks last? (\textbf{The long vertical on the right}.)

Source: \href{https://raw.githubusercontent.com/chanind/hanzi-writer-data/68d10a4b21150cae5e1ebbd223eed289cf32d90c/data/\%E5\%89\%8D.json}{Hanzi Writer Data \zh{前}, pinned stroke paths and medians}.
\end{tcolorbox}
\section[niánjì]{niánjì — age (of a person)}

\label{lesson:ZH-C78-nianji}



\begin{quote}
Before the new one: say the Mandarin for to grow up, then the Mandarin for ago.
\end{quote}

\subsection*{You'll want to know: niánjì}
\emph{niánjì} — \textquotedblleft{}age (of a person)\textquotedblright{}.

\begin{quote}
\emph{nǐ bàba duō dà niánjì} — How old is your father?
\end{quote}

Its characters come later; for now it is a word for the ear.

\begin{cousinweb}
\emph{niánjì} is built from \zh{年} 'year' and jì, both inherited from Middle Chinese, \emph{nen} and \emph{kiX} (Guangyun, in Baxter's transcription); the dictionary glosses the second as a span of twelve years.
\end{cousinweb}

\begin{grammarlens}[title={What you've built}]
One more word about people.
\end{grammarlens}

\subsection*{Your turn}
\begin{itemize}
\raggedright
  \item \emph{Say it:} \emph{niánjì}
  \item \emph{Say it:} \emph{niánjì}, once more
  \item \emph{Say it:} \emph{nǐ bàba duō dà niánjì} — How old is your father?
  \item \emph{Recall:} say the Mandarin for to grow up, then the Mandarin for ago, then say \emph{niánjì} again
\end{itemize}

\begin{tcolorbox}[breakable,colback=teal!4,colframe=teal!35!black,title={Before you move on}]
What does niánjì mean? (\textquotedblleft{}Age (of a person)\textquotedblright{}.) Say it once more.
\end{tcolorbox}
\section[xiǎoshíhou]{xiǎoshíhou — in childhood; when one was little}

\label{lesson:ZH-C78-xiaoshihou}



\begin{quote}
Before the new one: say the Mandarin for ago, then the Mandarin for age.
\end{quote}

\subsection*{You'll want to know: xiǎoshíhou}
\emph{xiǎoshíhou} — \textquotedblleft{}in childhood; when one was little\textquotedblright{}.

\begin{quote}
\emph{wǒ xiǎoshíhou hěn xǐhuan xiào} — As a child I laughed a lot.
\end{quote}

Its characters come later; for now it is a word for the ear.

\begin{cousinweb}
\emph{xiǎoshíhou} is built from \zh{小} 'small' and shíhou 'time', all inherited from Middle Chinese, \emph{sjewX}, \emph{dzyi} and \emph{huwH} (Guangyun, in Baxter's transcription): the time when you were small.
\end{cousinweb}

\begin{grammarlens}[title={What you've built}]
One more word about people.
\end{grammarlens}

\subsection*{Your turn}
\begin{itemize}
\raggedright
  \item \emph{Say it:} \emph{xiǎoshíhou}
  \item \emph{Say it:} \emph{xiǎoshíhou}, once more
  \item \emph{Say it:} \emph{wǒ xiǎoshíhou hěn xǐhuan xiào} — As a child I laughed a lot.
  \item \emph{Recall:} say the Mandarin for ago, then the Mandarin for age, then say \emph{xiǎoshíhou} again
\end{itemize}

\begin{tcolorbox}[breakable,colback=teal!4,colframe=teal!35!black,title={Before you move on}]
What does xiǎoshíhou mean? (\textquotedblleft{}In childhood; when one was little\textquotedblright{}.) Say it once more.
\end{tcolorbox}
